\section{Chapter Preface}
Computation can be viewed as motion rather than instruction execution: trajectories move through representational space toward progressively lower incompatibility. This dynamical lens is central to CLIO.

Seen this way, iterative reconstruction is a stability problem. We ask whether updates contract error, whether admissible sets are invariant, and whether oscillation can be controlled or characterized.

This chapter supplies the dynamical core by connecting reconstruction maps to Lyapunov and invariance arguments.

\section{Derivations and Formal Results}
Define iteration
\[
x_{t+1}=F(x_t;y,\constraintSet)
\]
and admissible target set
\[
A_y=\mathcal F_y\cap A.
\]
Define error \(E_t=d(x_t,A_y)\).

\begin{theorem}[Geometric error decay]
If \(E_{t+1}\le\rho E_t\) with \(0\le\rho<1\), then
\[
E_t\le\rho^tE_0\to0.
\]
\end{theorem}

\begin{proof}
Induction gives \(E_t\le\rho^tE_0\). Since \(\rho^t\to0\), convergence follows.
\end{proof}

\begin{definition}[Lyapunov witness for admissible convergence]
A function \(V:\stateSpace\to\mathbb R_{\ge0}\) is admissibility Lyapunov if
\[
V(x)=0\iff x\in A_y,
\quad
V(F(x))-V(x)\le0.
\]
\end{definition}

\begin{theorem}[Stability and asymptotic convergence]
If \(V\) is nonincreasing along trajectories, \(A_y\) is stable. If strict decrease holds for \(x\notin A_y\), trajectories converge asymptotically to \(A_y\).
\end{theorem}
